% Blunder hazard: the paper. What is typed here says what the numbers mean.
% Every number, table, figure and board is written by
% `lab paper blunder-hazard` into generated/, from the model's committed
% results.
\documentclass[11pt]{article}
% The preamble every paper of the lab shares: % The preamble every paper of the lab shares: % The preamble every paper of the lab shares: \input{../../../papers/preamble}
% from models/<model>/paper/, \input{../preamble} from papers/<name>/.
% Boards come from chessboard and xskak, plots from pgfplots. A paper prints no
% number by hand: `lab paper` writes every number, table and figure under the
% paper's generated/ folder, and paper.tex says what they mean.
\usepackage[LSBC4,T1]{fontenc}
\usepackage[utf8]{inputenc}
\usepackage{lmodern}
\usepackage[margin=1in]{geometry}
\usepackage{microtype}
\usepackage{amsmath}
\usepackage{booktabs}
\usepackage{tabularx}
\usepackage{array}
\usepackage{xcolor}
\usepackage[obeyspaces,spaces]{url}
\usepackage{xskak}
\usepackage{pgfplots}
\usepgfplotslibrary{groupplots}
\usetikzlibrary{positioning,arrows.meta,calc,fit}
\pgfplotsset{compat=1.18}
\usepackage{caption}
\usepackage{subcaption}
\usepackage{enumitem}
\usepackage[section]{placeins}
\usepackage[hidelinks]{hyperref}

% A generated value by its key: generated/values.tex sets them, \V{key} prints
% one, and a key nobody set stops the build.
\makeatletter
\newcommand{\setvalue}[2]{\expandafter\gdef\csname labvalue@#1\endcsname{#2}}
\newcommand{\V}[1]{\@ifundefined{labvalue@#1}%
  {\PackageError{lab}{No generated value named '#1'}{Run `lab paper` for this model again.}}%
  {\csname labvalue@#1\endcsname}}
\makeatother

% Chart colours. A run keeps its colour in every figure of every paper.
\definecolor{labBlue}{HTML}{2A78D6}
\definecolor{labOrange}{HTML}{EB6834}
\definecolor{labAqua}{HTML}{1BAF7A}
\definecolor{labYellow}{HTML}{EDA100}
\definecolor{labMagenta}{HTML}{E87BA4}
\definecolor{labGreen}{HTML}{008300}
\definecolor{labViolet}{HTML}{6250D6}
\definecolor{labGray}{HTML}{898781}
\definecolor{labInk}{HTML}{0B0B0B}
\definecolor{labSecondary}{HTML}{52514E}
\definecolor{labGrid}{HTML}{E1E0D9}
\definecolor{labAxis}{HTML}{C3C2B7}
\definecolor{labWash}{HTML}{CDE2FB}

\pgfplotsset{
  lab/.style={
    font=\sffamily\footnotesize,
    axis line style={labAxis},
    tick style={labAxis},
    tick label style={labSecondary},
    label style={labSecondary},
    title style={labInk, font=\sffamily\small\bfseries, align=left},
    grid style={labGrid, line width=0.3pt},
    major grid style={labGrid, line width=0.3pt},
    axis lines*=left,
    tick align=outside,
    clip=false,
    legend style={draw=none, fill=none, font=\sffamily\footnotesize, text=labSecondary,
      cells={anchor=west}, legend columns=-1, column sep=6pt, at={(0,1.03)}, anchor=south west},
    every axis plot/.append style={line width=1.1pt, line join=round, line cap=round},
    scaled ticks=false,
    /pgf/number format/1000 sep={,},
  },
  lab marks/.style={mark=*, mark size=1.9pt, mark options={draw=white, line width=0.6pt}},
}

\setchessboard{boardfontencoding=LSBC4, boardfontfamily=skaknew, blackfieldcolor=black!18,
  whitefieldcolor=white, setfontcolors, showmover=false, label=true, labelfontsize=6pt,
  boardfontsize=12pt}

% A move or line in UCI, a position in FEN (breakable at a slash), and a run's name.
\newcommand{\uci}[1]{\texttt{#1}}
\newcommand{\fen}[1]{{\footnotesize\path{#1}}}
\newcommand{\run}[1]{\texttt{#1}}
\newcolumntype{L}{>{\raggedright\arraybackslash}X}
\newcolumntype{R}{>{\raggedleft\arraybackslash}X}

\setlength{\parskip}{0.4em}
\setlength{\parindent}{0pt}
\captionsetup{font=small, labelfont=bf}

% from models/<model>/paper/, \input{../preamble} from papers/<name>/.
% Boards come from chessboard and xskak, plots from pgfplots. A paper prints no
% number by hand: `lab paper` writes every number, table and figure under the
% paper's generated/ folder, and paper.tex says what they mean.
\usepackage[LSBC4,T1]{fontenc}
\usepackage[utf8]{inputenc}
\usepackage{lmodern}
\usepackage[margin=1in]{geometry}
\usepackage{microtype}
\usepackage{amsmath}
\usepackage{booktabs}
\usepackage{tabularx}
\usepackage{array}
\usepackage{xcolor}
\usepackage[obeyspaces,spaces]{url}
\usepackage{xskak}
\usepackage{pgfplots}
\usepgfplotslibrary{groupplots}
\usetikzlibrary{positioning,arrows.meta,calc,fit}
\pgfplotsset{compat=1.18}
\usepackage{caption}
\usepackage{subcaption}
\usepackage{enumitem}
\usepackage[section]{placeins}
\usepackage[hidelinks]{hyperref}

% A generated value by its key: generated/values.tex sets them, \V{key} prints
% one, and a key nobody set stops the build.
\makeatletter
\newcommand{\setvalue}[2]{\expandafter\gdef\csname labvalue@#1\endcsname{#2}}
\newcommand{\V}[1]{\@ifundefined{labvalue@#1}%
  {\PackageError{lab}{No generated value named '#1'}{Run `lab paper` for this model again.}}%
  {\csname labvalue@#1\endcsname}}
\makeatother

% Chart colours. A run keeps its colour in every figure of every paper.
\definecolor{labBlue}{HTML}{2A78D6}
\definecolor{labOrange}{HTML}{EB6834}
\definecolor{labAqua}{HTML}{1BAF7A}
\definecolor{labYellow}{HTML}{EDA100}
\definecolor{labMagenta}{HTML}{E87BA4}
\definecolor{labGreen}{HTML}{008300}
\definecolor{labViolet}{HTML}{6250D6}
\definecolor{labGray}{HTML}{898781}
\definecolor{labInk}{HTML}{0B0B0B}
\definecolor{labSecondary}{HTML}{52514E}
\definecolor{labGrid}{HTML}{E1E0D9}
\definecolor{labAxis}{HTML}{C3C2B7}
\definecolor{labWash}{HTML}{CDE2FB}

\pgfplotsset{
  lab/.style={
    font=\sffamily\footnotesize,
    axis line style={labAxis},
    tick style={labAxis},
    tick label style={labSecondary},
    label style={labSecondary},
    title style={labInk, font=\sffamily\small\bfseries, align=left},
    grid style={labGrid, line width=0.3pt},
    major grid style={labGrid, line width=0.3pt},
    axis lines*=left,
    tick align=outside,
    clip=false,
    legend style={draw=none, fill=none, font=\sffamily\footnotesize, text=labSecondary,
      cells={anchor=west}, legend columns=-1, column sep=6pt, at={(0,1.03)}, anchor=south west},
    every axis plot/.append style={line width=1.1pt, line join=round, line cap=round},
    scaled ticks=false,
    /pgf/number format/1000 sep={,},
  },
  lab marks/.style={mark=*, mark size=1.9pt, mark options={draw=white, line width=0.6pt}},
}

\setchessboard{boardfontencoding=LSBC4, boardfontfamily=skaknew, blackfieldcolor=black!18,
  whitefieldcolor=white, setfontcolors, showmover=false, label=true, labelfontsize=6pt,
  boardfontsize=12pt}

% A move or line in UCI, a position in FEN (breakable at a slash), and a run's name.
\newcommand{\uci}[1]{\texttt{#1}}
\newcommand{\fen}[1]{{\footnotesize\path{#1}}}
\newcommand{\run}[1]{\texttt{#1}}
\newcolumntype{L}{>{\raggedright\arraybackslash}X}
\newcolumntype{R}{>{\raggedleft\arraybackslash}X}

\setlength{\parskip}{0.4em}
\setlength{\parindent}{0pt}
\captionsetup{font=small, labelfont=bf}

% from models/<model>/paper/, \input{../preamble} from papers/<name>/.
% Boards come from chessboard and xskak, plots from pgfplots. A paper prints no
% number by hand: `lab paper` writes every number, table and figure under the
% paper's generated/ folder, and paper.tex says what they mean.
\usepackage[LSBC4,T1]{fontenc}
\usepackage[utf8]{inputenc}
\usepackage{lmodern}
\usepackage[margin=1in]{geometry}
\usepackage{microtype}
\usepackage{amsmath}
\usepackage{booktabs}
\usepackage{tabularx}
\usepackage{array}
\usepackage{xcolor}
\usepackage[obeyspaces,spaces]{url}
\usepackage{xskak}
\usepackage{pgfplots}
\usepgfplotslibrary{groupplots}
\usetikzlibrary{positioning,arrows.meta,calc,fit}
\pgfplotsset{compat=1.18}
\usepackage{caption}
\usepackage{subcaption}
\usepackage{enumitem}
\usepackage[section]{placeins}
\usepackage[hidelinks]{hyperref}

% A generated value by its key: generated/values.tex sets them, \V{key} prints
% one, and a key nobody set stops the build.
\makeatletter
\newcommand{\setvalue}[2]{\expandafter\gdef\csname labvalue@#1\endcsname{#2}}
\newcommand{\V}[1]{\@ifundefined{labvalue@#1}%
  {\PackageError{lab}{No generated value named '#1'}{Run `lab paper` for this model again.}}%
  {\csname labvalue@#1\endcsname}}
\makeatother

% Chart colours. A run keeps its colour in every figure of every paper.
\definecolor{labBlue}{HTML}{2A78D6}
\definecolor{labOrange}{HTML}{EB6834}
\definecolor{labAqua}{HTML}{1BAF7A}
\definecolor{labYellow}{HTML}{EDA100}
\definecolor{labMagenta}{HTML}{E87BA4}
\definecolor{labGreen}{HTML}{008300}
\definecolor{labViolet}{HTML}{6250D6}
\definecolor{labGray}{HTML}{898781}
\definecolor{labInk}{HTML}{0B0B0B}
\definecolor{labSecondary}{HTML}{52514E}
\definecolor{labGrid}{HTML}{E1E0D9}
\definecolor{labAxis}{HTML}{C3C2B7}
\definecolor{labWash}{HTML}{CDE2FB}

\pgfplotsset{
  lab/.style={
    font=\sffamily\footnotesize,
    axis line style={labAxis},
    tick style={labAxis},
    tick label style={labSecondary},
    label style={labSecondary},
    title style={labInk, font=\sffamily\small\bfseries, align=left},
    grid style={labGrid, line width=0.3pt},
    major grid style={labGrid, line width=0.3pt},
    axis lines*=left,
    tick align=outside,
    clip=false,
    legend style={draw=none, fill=none, font=\sffamily\footnotesize, text=labSecondary,
      cells={anchor=west}, legend columns=-1, column sep=6pt, at={(0,1.03)}, anchor=south west},
    every axis plot/.append style={line width=1.1pt, line join=round, line cap=round},
    scaled ticks=false,
    /pgf/number format/1000 sep={,},
  },
  lab marks/.style={mark=*, mark size=1.9pt, mark options={draw=white, line width=0.6pt}},
}

\setchessboard{boardfontencoding=LSBC4, boardfontfamily=skaknew, blackfieldcolor=black!18,
  whitefieldcolor=white, setfontcolors, showmover=false, label=true, labelfontsize=6pt,
  boardfontsize=12pt}

% A move or line in UCI, a position in FEN (breakable at a slash), and a run's name.
\newcommand{\uci}[1]{\texttt{#1}}
\newcommand{\fen}[1]{{\footnotesize\path{#1}}}
\newcommand{\run}[1]{\texttt{#1}}
\newcolumntype{L}{>{\raggedright\arraybackslash}X}
\newcolumntype{R}{>{\raggedleft\arraybackslash}X}

\setlength{\parskip}{0.4em}
\setlength{\parindent}{0pt}
\captionsetup{font=small, labelfont=bf}

% Written by `lab paper blunder-hazard`; do not edit by hand.
\looktrue
\chancetrue
\sequencefalse
\contextgbdttrue
\pairstrue
\variantstrue
\shapestrue
\gamestrue
\featurestrue
\changestrue
\changehandtrue
\changeepochstrue
\changestarttrue
\changeattackstrue
\changegeometrytrue
\changeconvtrue
\changehybridtrue
\changelargerfalse
\moretrue
\shownfalse
\shownpairedfalse
\setvalue{band.1.floor}{0.1066}
\setvalue{band.1.name}{\textless{}1100}
\setvalue{band.5.floor}{0.0380}
\setvalue{band.5.name}{2000+}
\setvalue{bands.best-ahead.of}{all five bands}
\setvalue{bands.count}{five}
\setvalue{bands.lead.narrowest}{+0.0183}
\setvalue{bands.lead.narrowest.name}{1400-1700}
\setvalue{bands.lead.widest}{+0.0373}
\setvalue{bands.lead.widest.name}{\textless{}1100}
\setvalue{best.beside.closer}{48}
\setvalue{best.beside.correlation}{0.80}
\setvalue{best.bottom.50.rate}{99.2}
\setvalue{best.brier}{0.0558}
\setvalue{best.calibration-error}{0.0023}
\setvalue{best.calibration.groups}{ten}
\setvalue{best.calibration.highest}{27.6}
\setvalue{best.calibration.lowest}{0.0}
\setvalue{best.calibration.widest-gap}{0.9}
\setvalue{best.device}{NVIDIA A100-SXM4-40GB}
\setvalue{best.epochs.run}{12}
\setvalue{best.gain.of-lift}{13}
\setvalue{best.gain.percent}{9.6}
\setvalue{best.log-loss}{0.1959}
\setvalue{best.name}{\run{transformer-medium-hand-pretrained}}
\setvalue{best.parameters}{6{,}557{,}505}
\setvalue{best.platt.slope}{0.82}
\setvalue{best.roc-auc}{0.8403}
\setvalue{best.rows-a-second}{14{,}000}
\setvalue{best.score}{0.2776}
\setvalue{best.score.high}{0.2921}
\setvalue{best.score.low}{0.2643}
\setvalue{best.seconds}{9 min}
\setvalue{best.top.1.lift}{7.1}
\setvalue{best.top.1.rate}{48.8}
\setvalue{best.top.1.recall}{7.1}
\setvalue{best.top.10.rate}{27.5}
\setvalue{best.top.10.recall}{40.2}
\setvalue{best.top.20.rate}{21.7}
\setvalue{best.top.20.recall}{63.4}
\setvalue{best.top.50.recall}{94.1}
\setvalue{challenger.name}{\run{transformer-medium-hand-pretrained}}
\setvalue{code.commit}{\texttt{36a0ad9}}
\setvalue{context.lift.share}{65}
\setvalue{cost.seconds}{49 min}
\setvalue{cost.seconds.networks}{47 min}
\setvalue{cost.seconds.others}{3 min}
\setvalue{data.fingerprint}{\texttt{340c68ed5c91}}
\setvalue{data.games.test}{2{,}944}
\setvalue{data.licence}{CC0 1.0}
\setvalue{data.month}{October 2026}
\setvalue{data.name}{\texttt{game-positions-2026-10}}
\setvalue{data.rows}{698{,}013}
\setvalue{data.rows.test}{68{,}579}
\setvalue{data.share.test}{10}
\setvalue{data.share.train}{85}
\setvalue{data.share.valid}{5}
\setvalue{examples.count}{three}
\setvalue{floor.brier}{0.0639}
\setvalue{floor.constant}{6.7\%}
\setvalue{floor.log-loss}{0.2499}
\setvalue{floor.score}{0.0685}
\setvalue{grid.rating-clock.-1100.-30.rows}{47}
\setvalue{grid.rating-clock.1700-2000.-30.rows}{223}
\setvalue{grid.rating-win.-1100.60-80.said}{17.6\%}
\setvalue{grid.rating-win.2000+.60-80.happened}{5.9\%}
\setvalue{grid.rating-win.2000+.60-80.said}{6.9\%}
\setvalue{grid.rating-win.gbm-hand.widest-off}{1.8}
\setvalue{grid.rating-win.happened.most}{20.4\%}
\setvalue{grid.rating-win.transformer-medium-hand-pretrained.widest-off}{3.1}
\setvalue{grid.rating-win.transformer-medium-pretrained.widest-off}{2.3}
\setvalue{grid.smallest-cell}{30}
\setvalue{label.rate}{6.7}
\setvalue{label.rate.test}{6.9}
\setvalue{label.rate.train}{6.7}
\setvalue{label.rate.valid}{6.5}
\setvalue{ladder.behind.of}{three of the four networks}
\setvalue{look.date}{23 November 2026}
\setvalue{look.resamples}{200}
\setvalue{look.smallest-compared}{200}
\setvalue{look.test-rows}{68{,}579}
\setvalue{manifest.games-read}{1{,}003{,}668}
\setvalue{method.blunder-bar}{20}
\setvalue{method.gbm.features.percent}{80}
\setvalue{method.gbm.leaves}{127}
\setvalue{method.gbm.patience}{100}
\setvalue{method.gbm.rate}{0.05}
\setvalue{method.gbm.rounds}{4{,}000}
\setvalue{method.gbm.rows.percent}{80}
\setvalue{method.resamples}{1{,}000}
\setvalue{more.gbm-hand.most}{four times}
\setvalue{networks.ahead.of}{four of the seven networks}
\setvalue{networks.patience}{three}
\setvalue{networks.warmup.percent}{5}
\setvalue{pair.conv-medium-hand.transformer-medium-hand.gain}{$-$0.000209}
\setvalue{pair.conv-medium-hand.transformer-medium-hand.gain.high}{+0.005750}
\setvalue{pair.conv-medium-hand.transformer-medium-hand.gain.low}{$-$0.005854}
\setvalue{pair.conv-medium-hand.transformer-medium-hand.kept}{not accepted}
\setvalue{pair.conv-medium-hand.transformer-medium-hand.valid.gain}{+0.00608}
\setvalue{pair.conv-medium-hand.transformer-medium-hand.verdict}{not significantly different}
\setvalue{pair.gbm-hand-double.gbm-hand.gain}{+0.01017}
\setvalue{pair.gbm-hand-double.gbm-hand.kept}{accepted}
\setvalue{pair.gbm-hand-double.gbm-hand.verdict}{significantly better}
\setvalue{pair.gbm-hand-half.gbm-hand.gain}{$-$0.004299}
\setvalue{pair.gbm-hand-half.gbm-hand.verdict}{not significantly different}
\setvalue{pair.gbm-hand-no-hanging.gbm-hand.gain}{$-$0.00880}
\setvalue{pair.gbm-hand-no-kings.gbm-hand.gain}{$-$0.01014}
\setvalue{pair.gbm-hand-no-kings.gbm-hand.verdict}{significantly worse}
\setvalue{pair.gbm-hand-no-mobility.gbm-hand.gain}{$-$0.005001}
\setvalue{pair.gbm-hand-no-pawns.gbm-hand.gain}{$-$0.000488}
\setvalue{pair.gbm-hand-no-pawns.gbm-hand.verdict}{not significantly different}
\setvalue{pair.gbm-hand-quadruple.gbm-hand.gain}{+0.01256}
\setvalue{pair.gbm-hand-quadruple.gbm-hand.kept}{accepted}
\setvalue{pair.gbm-hand-quadruple.gbm-hand.verdict}{significantly better}
\setvalue{pair.gbm-hand-quarter.gbm-hand.gain}{$-$0.01203}
\setvalue{pair.gbm-hand-quarter.gbm-hand.verdict}{significantly worse}
\setvalue{pair.hybrid-medium-hand.transformer-medium-hand.gain}{$-$0.00438}
\setvalue{pair.hybrid-medium-hand.transformer-medium-hand.gain.high}{+0.00153}
\setvalue{pair.hybrid-medium-hand.transformer-medium-hand.gain.low}{$-$0.01037}
\setvalue{pair.hybrid-medium-hand.transformer-medium-hand.kept}{not accepted}
\setvalue{pair.hybrid-medium-hand.transformer-medium-hand.verdict}{not significantly different}
\setvalue{pair.transformer-medium-attacks.transformer-medium.gain}{+0.001142}
\setvalue{pair.transformer-medium-attacks.transformer-medium.gain.high}{+0.004468}
\setvalue{pair.transformer-medium-attacks.transformer-medium.gain.low}{$-$0.001737}
\setvalue{pair.transformer-medium-attacks.transformer-medium.kept}{not accepted}
\setvalue{pair.transformer-medium-attacks.transformer-medium.valid.gain}{+0.000815}
\setvalue{pair.transformer-medium-attacks.transformer-medium.verdict}{not significantly different}
\setvalue{pair.transformer-medium-hand-attacks.transformer-medium-hand.gain}{$-$0.001845}
\setvalue{pair.transformer-medium-hand-attacks.transformer-medium-hand.gain.high}{+0.000665}
\setvalue{pair.transformer-medium-hand-attacks.transformer-medium-hand.gain.low}{$-$0.004391}
\setvalue{pair.transformer-medium-hand-attacks.transformer-medium-hand.kept}{not accepted}
\setvalue{pair.transformer-medium-hand-attacks.transformer-medium-hand.verdict}{not significantly different}
\setvalue{pair.transformer-medium-hand-double.transformer-medium-hand.gain}{+0.001259}
\setvalue{pair.transformer-medium-hand-double.transformer-medium-hand.kept}{not accepted}
\setvalue{pair.transformer-medium-hand-double.transformer-medium-hand.verdict}{not significantly different}
\setvalue{pair.transformer-medium-hand-half.transformer-medium-hand.gain}{$-$0.00848}
\setvalue{pair.transformer-medium-hand-long.transformer-medium-hand.gain}{$-$0.003788}
\setvalue{pair.transformer-medium-hand-long.transformer-medium-hand.gain.high}{+0.000907}
\setvalue{pair.transformer-medium-hand-long.transformer-medium-hand.gain.low}{$-$0.008039}
\setvalue{pair.transformer-medium-hand-long.transformer-medium-hand.kept}{not accepted}
\setvalue{pair.transformer-medium-hand-long.transformer-medium-hand.valid.verdict}{no better}
\setvalue{pair.transformer-medium-hand-long.transformer-medium-hand.verdict}{not significantly different}
\setvalue{pair.transformer-medium-hand-no-geometry.transformer-medium-hand.gain}{$-$0.001523}
\setvalue{pair.transformer-medium-hand-no-geometry.transformer-medium-hand.gain.high}{+0.001176}
\setvalue{pair.transformer-medium-hand-no-geometry.transformer-medium-hand.gain.low}{$-$0.004325}
\setvalue{pair.transformer-medium-hand-no-geometry.transformer-medium-hand.verdict}{not significantly different}
\setvalue{pair.transformer-medium-hand-no-hanging.transformer-medium-hand.gain}{$-$0.01019}
\setvalue{pair.transformer-medium-hand-no-kings.transformer-medium-hand.gain}{+0.000052}
\setvalue{pair.transformer-medium-hand-no-kings.transformer-medium-hand.verdict}{not significantly different}
\setvalue{pair.transformer-medium-hand-no-mobility.transformer-medium-hand.gain}{$-$0.00948}
\setvalue{pair.transformer-medium-hand-no-pawns.transformer-medium-hand.gain}{+0.001076}
\setvalue{pair.transformer-medium-hand-no-pawns.transformer-medium-hand.verdict}{not significantly different}
\setvalue{pair.transformer-medium-hand-pretrained.transformer-medium-hand.gain}{+0.01098}
\setvalue{pair.transformer-medium-hand-pretrained.transformer-medium-hand.gain.high}{+0.01966}
\setvalue{pair.transformer-medium-hand-pretrained.transformer-medium-hand.gain.low}{+0.00234}
\setvalue{pair.transformer-medium-hand-pretrained.transformer-medium-hand.kept}{accepted}
\setvalue{pair.transformer-medium-hand-pretrained.transformer-medium-hand.valid.verdict}{better}
\setvalue{pair.transformer-medium-hand-pretrained.transformer-medium-hand.verdict}{significantly better}
\setvalue{pair.transformer-medium-hand-quadruple.transformer-medium-hand.gain}{+0.004459}
\setvalue{pair.transformer-medium-hand-quadruple.transformer-medium-hand.kept}{accepted}
\setvalue{pair.transformer-medium-hand-quadruple.transformer-medium-hand.verdict}{significantly better}
\setvalue{pair.transformer-medium-hand-quarter.transformer-medium-hand.gain}{$-$0.02531}
\setvalue{pair.transformer-medium-hand.transformer-medium.gain}{+0.02962}
\setvalue{pair.transformer-medium-hand.transformer-medium.gain.high}{+0.03802}
\setvalue{pair.transformer-medium-hand.transformer-medium.gain.low}{+0.02098}
\setvalue{pair.transformer-medium-hand.transformer-medium.kept}{accepted}
\setvalue{pair.transformer-medium-hand.transformer-medium.valid.verdict}{better}
\setvalue{pair.transformer-medium-hand.transformer-medium.verdict}{significantly better}
\setvalue{pair.transformer-medium-long.transformer-medium.gain}{$-$0.001048}
\setvalue{pair.transformer-medium-long.transformer-medium.gain.high}{+0.002857}
\setvalue{pair.transformer-medium-long.transformer-medium.gain.low}{$-$0.005152}
\setvalue{pair.transformer-medium-long.transformer-medium.kept}{not accepted}
\setvalue{pair.transformer-medium-long.transformer-medium.valid.verdict}{better}
\setvalue{pair.transformer-medium-long.transformer-medium.verdict}{not significantly different}
\setvalue{pair.transformer-medium-pretrained.transformer-medium.gain}{+0.03773}
\setvalue{pair.transformer-medium-pretrained.transformer-medium.gain.high}{+0.04832}
\setvalue{pair.transformer-medium-pretrained.transformer-medium.gain.low}{+0.02816}
\setvalue{pair.transformer-medium-pretrained.transformer-medium.kept}{accepted}
\setvalue{pair.transformer-medium-pretrained.transformer-medium.valid.verdict}{better}
\setvalue{pair.transformer-medium-pretrained.transformer-medium.verdict}{significantly better}
\setvalue{recipe.games}{30{,}000}
\setvalue{recipe.min-ply}{11}
\setvalue{recipe.per-game}{24}
\setvalue{reference.brier}{0.0569}
\setvalue{reference.calibration-error}{0.0042}
\setvalue{reference.calibration.widest-gap}{0.9}
\setvalue{reference.features}{25}
\setvalue{reference.leans.1.feature}{win probability before the move}
\setvalue{reference.leans.1.share}{41}
\setvalue{reference.leans.2.feature}{material balance}
\setvalue{reference.leans.2.share}{11}
\setvalue{reference.leans.3.feature}{mover's rating}
\setvalue{reference.leans.3.share}{7}
\setvalue{reference.log-loss}{0.2013}
\setvalue{reference.platt.slope}{1.03}
\setvalue{reference.roc-auc}{0.8232}
\setvalue{reference.score}{0.2533}
\setvalue{reference.trees}{196}
\setvalue{results.date}{23 November 2026}
\setvalue{run.conv-medium-hand.lr}{0.002}
\setvalue{run.conv-medium-hand.parameters.millions}{6.9}
\setvalue{run.conv-small.best-epoch}{7}
\setvalue{run.conv-small.epochs}{8}
\setvalue{run.conv-small.gain}{$-$0.01642}
\setvalue{run.conv-small.gain.high}{$-$0.00713}
\setvalue{run.conv-small.gain.low}{$-$0.02564}
\setvalue{run.conv-small.parameters.millions}{1.0}
\setvalue{run.conv-small.score}{0.2369}
\setvalue{run.gbm-context.roc-auc}{0.7797}
\setvalue{run.gbm-context.score}{0.1879}
\setvalue{run.gbm-hand-quadruple.score}{0.2658}
\setvalue{run.gbm-hand-quadruple.trees}{503}
\setvalue{run.gbm-hand-quarter.score}{0.2413}
\setvalue{run.hybrid-medium-hand.convs}{2}
\setvalue{run.hybrid-medium-hand.depth}{5}
\setvalue{run.hybrid-medium-hand.parameters.millions}{6.6}
\setvalue{run.transformer-medium-attacks.gain}{$-$0.01517}
\setvalue{run.transformer-medium-attacks.gain.high}{$-$0.00659}
\setvalue{run.transformer-medium-attacks.gain.low}{$-$0.02352}
\setvalue{run.transformer-medium-attacks.score}{0.2381}
\setvalue{run.transformer-medium-hand-double.best-epoch}{6}
\setvalue{run.transformer-medium-hand-double.epochs}{6}
\setvalue{run.transformer-medium-hand-long.best-epoch}{14}
\setvalue{run.transformer-medium-hand-pretrained.best-epoch}{9}
\setvalue{run.transformer-medium-hand-pretrained.epochs}{12}
\setvalue{run.transformer-medium-hand-pretrained.over-twin}{+0.00288}
\setvalue{run.transformer-medium-hand-pretrained.score}{0.2776}
\setvalue{run.transformer-medium-hand-quadruple.best-epoch}{3}
\setvalue{run.transformer-medium-hand-quadruple.epochs}{3}
\setvalue{run.transformer-medium-hand-quadruple.score}{0.2711}
\setvalue{run.transformer-medium-hand.best-epoch}{9}
\setvalue{run.transformer-medium-hand.epochs}{12}
\setvalue{run.transformer-medium-hand.gain}{+0.01331}
\setvalue{run.transformer-medium-hand.gain.high}{+0.01978}
\setvalue{run.transformer-medium-hand.gain.low}{+0.00701}
\setvalue{run.transformer-medium-hand.lr}{0.0006}
\setvalue{run.transformer-medium-hand.parameters.millions}{6.6}
\setvalue{run.transformer-medium-hand.platt.slope}{1.02}
\setvalue{run.transformer-medium-hand.score}{0.2666}
\setvalue{run.transformer-medium-long.best-epoch}{14}
\setvalue{run.transformer-medium-long.epochs}{24}
\setvalue{run.transformer-medium-long.gain}{$-$0.01736}
\setvalue{run.transformer-medium-long.gain.high}{$-$0.00891}
\setvalue{run.transformer-medium-long.gain.low}{$-$0.02569}
\setvalue{run.transformer-medium-pretrained.best-epoch}{10}
\setvalue{run.transformer-medium-pretrained.gain}{+0.02142}
\setvalue{run.transformer-medium-pretrained.gain.high}{+0.03066}
\setvalue{run.transformer-medium-pretrained.gain.low}{+0.01202}
\setvalue{run.transformer-medium-pretrained.log-loss}{0.1956}
\setvalue{run.transformer-medium-pretrained.platt.slope}{0.77}
\setvalue{run.transformer-medium-pretrained.roc-auc}{0.8411}
\setvalue{run.transformer-medium-pretrained.score}{0.2747}
\setvalue{run.transformer-medium.best-epoch}{12}
\setvalue{run.transformer-medium.epochs}{12}
\setvalue{run.transformer-medium.gain}{$-$0.01631}
\setvalue{run.transformer-medium.gain.high}{$-$0.00742}
\setvalue{run.transformer-medium.gain.low}{$-$0.02519}
\setvalue{run.transformer-medium.parameters.millions}{6.6}
\setvalue{run.transformer-medium.score}{0.2370}
\setvalue{run.transformer-small.best-epoch}{7}
\setvalue{run.transformer-small.epochs}{8}
\setvalue{run.transformer-small.gain}{$-$0.01692}
\setvalue{run.transformer-small.gain.high}{$-$0.00821}
\setvalue{run.transformer-small.gain.low}{$-$0.02552}
\setvalue{run.transformer-small.parameters.millions}{0.9}
\setvalue{run.transformer-small.score}{0.2364}
\setvalue{runs.count}{twelve}
\setvalue{slice.band.-1100.challenger.off}{$-$1.5}
\setvalue{slice.band.-1100.reference.off}{$-$0.9}
\setvalue{slice.band.2000+.challenger.off}{+0.6}
\setvalue{slice.band.2000+.reference.off}{$-$0.3}
\setvalue{slice.band.ahead.groups}{\textless{}1100, 1100-1400 and 1700-2000}
\setvalue{slice.band.ahead.of}{three of the five groups}
\setvalue{slice.band.gain.most.group}{\textless{}1100}
\setvalue{slice.check.no.happened}{7.1\%}
\setvalue{slice.check.yes.happened}{3.9\%}
\setvalue{slice.clock.-30.happened}{6.0\%}
\setvalue{slice.clock.-30.rows}{718}
\setvalue{slice.clock.30-60.happened}{9.0\%}
\setvalue{slice.clock.30-60.rows}{777}
\setvalue{slice.clock.480+.share}{58}
\setvalue{slice.clock.60-120.happened}{4.7\%}
\setvalue{slice.clock.ahead.groups}{480+}
\setvalue{slice.clock.ahead.of}{one of the six groups}
\setvalue{slice.hanging.0.happened}{6.3\%}
\setvalue{slice.hanging.0.share}{60}
\setvalue{slice.hanging.3+.happened}{8.2\%}
\setvalue{slice.moves.-5.happened}{3.0\%}
\setvalue{slice.moves.45+.happened}{9.0\%}
\setvalue{slice.moves.ahead.groups}{25-35, 35-45 and 45+}
\setvalue{slice.moves.ahead.of}{three of the six groups}
\setvalue{slice.open.-20.share}{20}
\setvalue{slice.opponent-hanging.0.happened}{6.8\%}
\setvalue{slice.opponent-hanging.3+.happened}{7.9\%}
\setvalue{slice.pieces.-3.gain}{$-$0.01837}
\setvalue{slice.pieces.3-6.gain}{+0.01917}
\setvalue{slice.pieces.6-9.gain}{+0.01414}
\setvalue{slice.pieces.ahead.groups}{9-12 and 12+}
\setvalue{slice.ply.30-40.happened}{8.5\%}
\setvalue{slice.ply.60-80.gain}{+0.00562}
\setvalue{slice.ply.80+.gain}{$-$0.02244}
\setvalue{slice.ply.80+.happened}{4.0\%}
\setvalue{slice.ply.80+.share}{6}
\setvalue{slice.ply.ahead.groups}{\textless{}20, 20-30, 30-40 and 40-60}
\setvalue{slice.ply.ahead.of}{four of the six groups}
\setvalue{slice.seconds.-2.challenger.off}{+0.9}
\setvalue{slice.seconds.-2.happened}{3.1\%}
\setvalue{slice.seconds.-2.reference.off}{+1.3}
\setvalue{slice.seconds.40+.challenger.off}{$-$3.0}
\setvalue{slice.seconds.40+.happened}{11.6\%}
\setvalue{slice.seconds.40+.reference.off}{$-$3.3}
\setvalue{slice.seconds.5-10.happened}{6.8\%}
\setvalue{slice.win.-10.share}{14}
\setvalue{slice.win.10-20.share}{6}
\setvalue{slice.win.20-40.happened}{8.8\%}
\setvalue{slice.win.40-60.happened}{7.5\%}
\setvalue{slice.win.60-80.happened}{12.4\%}
\setvalue{slice.win.80-90.gain}{+0.01592}
\setvalue{slice.win.80-90.happened}{14.3\%}
\setvalue{slice.win.90+.gain}{$-$0.00773}
\setvalue{slice.win.90+.happened}{5.4\%}
\setvalue{slice.win.ahead.groups}{20-40, 40-60 and 60-80}
\setvalue{slice.win.ahead.of}{three of the five groups}
\setvalue{slice.win.challenger.widest-off}{0.9}
\setvalue{slice.win.reference.widest-off}{1.0}
\setvalue{study.features.gbm-hand.most.gain}{$-$0.03451}
\setvalue{study.features.gbm-hand.most.group}{material}
\setvalue{study.features.gbm-hand.worse.of}{four of the five ablations}
\setvalue{study.features.transformer-medium-hand.most.gain}{$-$0.02723}
\setvalue{study.features.transformer-medium-hand.worse.of}{three of the five ablations}
\setvalue{study.games.transformer-medium-hand.worse.of}{both ablations}
\setvalue{task.context.count}{twelve}
\setvalue{task.hand.count}{13}
\setvalue{task.marks.count}{three}
\setvalue{task.reference}{\run{gbm-hand}}


\begin{document}

\labtitle{Blunder hazard}
  {Before a move is played, how likely is the player to blunder here?}
  {Results of \V{results.date}\quad\textperiodcentered\quad dataset \V{data.name},
   build \V{data.fingerprint}\quad\textperiodcentered\quad code \V{code.commit}}

\begin{summary}
Some positions are quiet and some are cliffs, and which is which depends on
the player's level and clock as much as on the board. This paper trains models
that say, before a move is played, how likely the mover is to blunder: to lose
more than \V{method.blunder-bar} points of win probability with the move. The
data is \V{data.rows} positions from \V{recipe.games} rated rapid games on
Lichess that carry engine analysis and clocks, split by game. A blunder is
rare, \V{floor.constant} of the moves. Ranked by the area under the
precision-recall curve, the best of \V{runs.count} runs is a square-token
transformer that reads the board and \V{task.hand.count} hand-made counts. It
scores
\V{best.score} (95\% interval \V{best.score.low} to \V{best.score.high}),
against \V{reference.score} for gradient-boosted trees on the hand-made counts
and \V{floor.score} for the base rate. Beyond the noise,
\V{networks.ahead.of} beat the trees. From the board alone no network of the
ladder matches them, whatever its size. Started from a body that was first
trained on engine evaluations, the same transformer passes them from the
board alone, and the best run starts from that body too. After calibration
the stated chance
can be read as a frequency: none of the best run's
\V{best.calibration.groups} groups is further than
\V{best.calibration.widest-gap} points from what happened.
\iflook
Of the hundredth of the test positions the best run calls most dangerous,
\V{best.top.1.rate}\% were blundered, \V{best.top.1.lift} times the base
rate.
\fi
\end{summary}

\section{The question}

An engine can say after the fact that a move was a blunder. It cannot say
beforehand that this is a position where people blunder. A hanging piece in a
quiet middlegame and a hanging piece with ten seconds left are the same to an
engine, and very different to the person at the board.

The coaching app this lab works for reviews a player's games. A chance of
blundering, given the position, the ratings and the clocks, would let it do
two things it cannot do today. It could mark the positions where players at
the player's level usually go wrong, which are the ones worth slowing down in,
whether or not the player went wrong there this time. And it could tell a
blunder in a position where almost nobody blunders, a lapse, from one in a
position that trips most players of that level, a hard position.

Human error in chess has been modelled before. With endgame tablebases as a
benchmark of perfection, the difficulty of the position was found to predict
a blunder better than the player's skill or the time left~\cite{humanerror},
and blunder prediction has been studied on Lichess games alongside human-move
prediction~\cite{maia}. This paper asks the question in the form the app
needs it: one calibrated chance for one position, from inputs the app has for
any game, compared against a plain baseline on games no learner saw.

\begin{table}[htbp]\centering\small
\begin{tabularx}{\textwidth}{lL}
\toprule
Reads & The position, the opponent's last move, both ratings, both clocks, the
time control, the move number, and the mover's win probability from the
engine's evaluation before the move. \\
Predicts & The chance that the move played loses more than
\V{method.blunder-bar} points of win probability. \\
Data & Rated Lichess rapid games of \V{data.month} (\V{data.licence}) with
engine analysis and clocks: up to \V{recipe.per-game} positions a game. \\
Split & By game: \V{data.share.train}\% train, \V{data.share.valid}\% valid,
\V{data.share.test}\% test. \\
Score & PR-AUC, the area under the precision-recall curve, on the test
positions. Higher is better. \\
\bottomrule
\end{tabularx}
\caption{The model in brief.}
\label{tab:brief}
\end{table}

\section{The data}

The source is the Lichess database of rated games for
\V{data.month}~\cite{lichessdb}. The builder reads the month as a stream from
its start and stops once \V{recipe.games} games have given rows, so the
sample is the first games of the month and not a draw from all of it. A game
gives rows when it is rated, played at rapid speed between two rated humans,
ended over the board or on time, and carries Lichess's engine analysis and
both clocks.

From each such game up to \V{recipe.per-game} positions are kept, chosen by a
hash of the game and the ply. None comes before ply \V{recipe.min-ply}, each
side's sixth move, while players are still playing their openings from
memory. A position is kept only when the evaluation before and after the move
and the clocks around it are all known, so every row carries every label.
That gives \V{data.rows} positions (Table~\ref{tab:data}). The split is a
hash of the game, so no game has positions on both sides of it.

\paragraph{What counts as a blunder.}
An engine's evaluation is turned into the mover's chance of winning by the
curve Lichess publishes~\cite{lichessaccuracy}, which the app uses as well; a
parity check in this repository holds the lab's curve to the app's. A move is
a blunder when it loses more than \V{method.blunder-bar} points of that
chance, the bar the app and Lichess share. By that bar \V{floor.constant} of
the moves in the train rows are blunders.

\begin{table}[htbp]\centering\small
\block{data}
\caption{The dataset's rows by split.}
\label{tab:data}
\end{table}

\iflook
The builder read \V{manifest.games-read} games of the month to find the
\V{manifest.games-kept} that gave rows. Over every row, \V{label.rate}\% of
the moves are blunders. Table~\ref{tab:describe} gives the spread of the
ratings, the clocks and a few other columns.

\begin{table}[htbp]\centering\small
\block{describe}
\caption{The spread of a few columns over every row of the dataset: the mean,
five quantiles, and the least and the most. Clocks and times are in seconds,
winning chances in percent.}
\label{tab:describe}
\end{table}
\fi

\section{What a model reads}
\label{sec:reads}

% How a position reaches a learner: the same for every model, so each paper
% reads it from here, under its own section heading, and adds what is its own.
Every model of the lab reads a position the same way, so two models differ in
what they are asked and never in how a board reaches them.

\paragraph{The board.}
A position arrives as a FEN and becomes 64 tokens, one a square. The mover's
pawn, knight, bishop, rook, queen and king are tokens 1 to 6, the opponent's
are 7 to 12, and an empty square is 0. When Black is to move the board is
turned and the colours are swapped, so the side to move always plays up the
board. A network then learns one game, not a white one and a black one, and a
position it meets with White to move also teaches it the mirrored position
with Black to move.

\paragraph{Marks.}
What the task knows about particular squares travels as marks: yes-or-no
channels laid over the 64 squares, which a network reads together with the
pieces. This model has \V{task.marks.count} (Table~\ref{tab:marks}). The first
is always the square a pawn may be taken on in passing. The others name the
squares a move leaves and reaches, which is how a network is shown a move
without being told what the move does.

\paragraph{Context numbers.}
Whatever is not about a square is a number beside the board: the four castling
rights, the mover's first, and what the task knows about the game. This model
has \V{task.context.count} (Table~\ref{tab:reads}). Every learner reads them,
the trees that see no board included. A network reads them standardised by the
mean and spread of the train rows.

\paragraph{Hand features.}
Gradient-boosted trees cannot read 64 tokens, so the lab also computes what a
person would count on a board: material, mobility, checks and captures on
offer, pieces that can be taken, passed pawns and pressure on each king, and
for a line of moves what the line does. Table~\ref{tab:hand} in the appendix
lists all \V{task.hand.count}. They serve twice. They are what the reference
run reads, which makes them the bar for a network: a network that reads the
squares themselves should find at least what a few dozen counts already hold.
And a run can hand them to a network beside the board, where they join the
context numbers, which tests whether the network had found them on its own.

\begin{table}[htbp]\centering\small
\block{reads}
\caption{The context numbers: what every learner reads beside the board.}
\label{tab:reads}
\end{table}

\begin{table}[htbp]\centering\small
\block{marks}
\caption{The marks: what the task tells a network about squares. A square
is always given as the mover sees it, after the board is turned.}
\label{tab:marks}
\end{table}

\iflook
\begin{figure}[htbp]\centering
\block{encoding}
\caption{One position of the dataset twice. Left: as it was played, with an
arrow for the opponent's move that led to it (grey) and for the move the row
is about (blue). Right: as a network reads it, turned so that the mover plays
up the board, with a wash on every square a mark is set on: grey for the
opponent's move and, where the task marks a line of moves, blue for the
mover's and orange for the replies.}
\label{fig:encoding}
\end{figure}
\fi


\paragraph{This model reads the position before the move, never the move.}
The marks carry only the squares of the opponent's last move. The context
numbers are both ratings, both clocks and the time control's starting time as
logarithms, the increment, the ply, and the mover's win probability before
the move. That last number says how the mover stands. It does not say which
move is best or how many moves keep the evaluation, so a learner knows that
the position is winning and not how easy the win is to spoil.

\section{The learners}
\label{sec:learners}

% The learners every model is given: the same ladder of runs on every task,
% so each paper reads it from here under its own section heading.
A run is one learner, trained on the train rows, tuned on the valid rows and
scored once on the test rows. Every model keeps the same ladder of runs, so
each step up it says what one more thing is worth. Table~\ref{tab:runs} lists
this model's.

\paragraph{The floor.}
\run{constant} says one number for every row: the mean label of the train
rows, or their median where the score is an absolute error. A model that does
not beat it has learned nothing.

\paragraph{Trees.}
\run{gbm-context} is gradient-boosted trees (LightGBM~\cite{lightgbm}) on the
context numbers alone, with no view of the board. It measures how much of the
answer is already in the numbers beside the board. \run{gbm-hand} adds the
hand features and is the \emph{reference}: every other run is compared with
it on the same test rows. Both grow up to \V{method.gbm.rounds} trees of up
to \V{method.gbm.leaves} leaves at a learning rate of \V{method.gbm.rate},
each tree on a random \V{method.gbm.rows.percent}\% of the rows and
\V{method.gbm.features.percent}\% of the features, and stop when the valid
rows have not improved for \V{method.gbm.patience} rounds.

\paragraph{Networks.}
A network is a body, which turns the tokens, the marks and the context numbers
into one vector, and a head of two layers, which turns that vector into the
answer. Two bodies ship, so that a leaderboard can say what the architecture
is worth on a task.

The \emph{square transformer}~\cite{attention} treats each square as a token:
an embedding of its piece, plus a learned embedding of the square, plus a
linear map of the square's marks. Every square attends to every other from the
first layer, so a bishop on a1 and a king on h8 meet at once. Attention is
told about the board's geometry by a learned bias for each offset between two
squares, one per head~\cite{relative}: files, ranks, diagonals and knight hops
are all offsets, so a head can learn to look along them. The context numbers
enter as one more token, and the body's output is that token beside the mean
of the 64 square tokens.

The \emph{convolutional tower} is residual blocks of three-by-three
convolutions over the board, the kind of network the first AlphaZero-style
engines used~\cite{alphazero}. The context numbers are added to every square
after the first layer, and the output is the mean and the maximum over the
squares. It sees further only as it gets deeper. It is cheap and trains on a
CPU, and it is here as the honest alternative: where it matches the
transformer, the cheaper model wins.

Sizes are settings of a run, not code. \run{conv-small} and
\run{transformer-small} are sized to train on a CPU. \run{transformer-medium}
is sized for a GPU, and \run{transformer-medium-hand} is the same network
reading the hand features beside the board. Table~\ref{tab:settings} in the
appendix has every size and schedule.

\paragraph{One training loop.}
Every network is trained the same way: AdamW~\cite{adamw}, a learning rate
that rises over the first \V{networks.warmup.percent}\% of the steps and then
follows a cosine down to zero, gradients clipped, mixed precision on a GPU.
The loss follows the score: squared error for a quantity ranked by RMSE, a
smoothed absolute error for one ranked by MAE, the log loss for a chance. A
quantity's label is standardised by the train rows. After each epoch the
network is scored on the valid rows, the weights of the best epoch are the
ones kept, and training stops early once \V{networks.patience} epochs pass
without a better one.

\begin{table}[htbp]\centering\footnotesize
\block{runs}
\caption{The runs, best first. \emph{Kept} is what the valid rows chose: the
number of trees, or the epoch whose weights were kept. \emph{Time} is the
whole run, training and scoring.}
\label{tab:runs}
\end{table}


\section{How a result is scored}
\label{sec:scoring}

The model's card named its score before the first full run: PR-AUC, the area
under the precision-recall curve, also called average precision. Put the test
positions in the order of the stated chance, highest first. At each position
that was a blunder, note the share of blunders among the positions so far.
PR-AUC is the average of those shares. A run that says the same chance
everywhere scores the base rate, and a run that puts every blunder first
scores one. It was chosen because a blunder is rare, and what the app needs
is the ranking: which positions to flag.

Four more scores are reported beside it. ROC-AUC is the chance that a blunder
is ranked above a move that was not one. The Brier score, the log loss and
the calibration error ask whether the stated chance means what it says, and
are computed after the chance is calibrated on the valid rows, so that a
stated one in five can be read as one in five. A run beats the reference when
its paired gain in PR-AUC has a 95\% interval above zero. Results are also
given by the mover's rating in \V{bands.count} bands.

% How every result of the lab is scored: the same for every model, so each
% paper names its own score under its section heading and reads the rest here.
\paragraph{Noise is measured over games.}
Rows that come from one game are not independent, so every interval in this
paper resamples whole games. The test games are drawn again with replacement
\V{method.resamples} times, the score is computed each time on the rows of the
games drawn, and the interval given is the middle 95\% of those scores.

\paragraph{A gain is paired.}
A run and the reference are scored on the same resampled games, and what is
reported is the difference, with its own interval. A gain counts when that
interval leaves zero out, and is then set in bold. Two runs that are both good
on easy games and bad on hard ones are compared fairly this way, where their
two separate intervals would overlap and hide a real difference.

\paragraph{Test rows are for the final score.}
\ifchance
The number of trees, the epoch whose weights are kept and the calibration of
the chance are all chosen on the valid rows.
\else
The number of trees and the epoch whose weights are kept are both chosen on
the valid rows.
\fi
The test rows are read once a run, for the scores printed here.

\ifchance
\paragraph{A chance is calibrated before it is scored.}
A model can rank well and still say four in ten where three in ten happen. So
one line is put through each run's log-odds, a slope and an intercept fitted
on the valid rows (Platt scaling~\cite{platt}), and applied to its test
predictions before they are scored. A run is then not marked down for a scale
that one line can fix, and what is left of its miscalibration is its own.
\fi

\paragraph{Reading the tables and figures.}
A run keeps one colour and one mark in every figure, and the reference is
always orange. Every figure has its numbers in a table, here or in the
appendix, so nothing has to be read off a colour. No number in this paper was
typed: each is written from the committed record of a run by the command that
builds the paper.


\section{Results}
\label{sec:results}

\subsection{The leaderboard}

Read from the bottom up, Table~\ref{tab:leaderboard} is a ladder. The floor
says the base rate for every position and scores \V{floor.score}. Trees on
the context numbers alone, \V{run.gbm-context.name}, reach
\V{run.gbm-context.score} without ever seeing a board: the ratings, the
clocks and above all the evaluation carry \V{context.lift.share}\% of what
the reference adds to the floor. The reference, which also reads the hand
features of the position, scores \V{reference.score}. The best run scores
\V{best.score} (Figure~\ref{fig:scores}).

\begin{table}[htbp]\centering\small
\block{leaderboard}
\caption{The runs on the \V{data.rows.test} test positions, from
\V{data.games.test} games, best first. The gain is paired with
\V{task.reference} on the same resampled games, positive when the run is
better, and bold when its interval leaves zero out.}
\label{tab:leaderboard}
\end{table}

\begin{figure}[htbp]\centering
\block{scores}
\caption{Every run's PR-AUC with its 95\% interval. The vertical line is the
reference's score.}
\label{fig:scores}
\end{figure}

\subsection{What a network adds}

Beyond the noise, \V{ladder.ahead.of} beat the trees and
\V{ladder.behind.of} fell short of them (Figure~\ref{fig:gains}). The one
ahead is the medium transformer that also reads the hand features, with a
paired gain of \V{run.transformer-medium-hand.gain}
(\V{run.transformer-medium-hand.gain.low} to
\V{run.transformer-medium-hand.gain.high}). In Figure~\ref{fig:scores} its
interval and the reference's overlap. With \V{data.games.test} test games,
only the paired comparison, which scores both runs on the same resampled
games, separates them.

\begin{figure}[htbp]\centering
\block{gains}
\caption{Each network's paired gain in PR-AUC over the reference, with its
95\% interval. Right of the line a network is ahead of the trees, left of it
behind.}
\label{fig:gains}
\end{figure}

From the board alone, no network of the ladder matches the trees. The
convolutional tower and both transformers without the hand features land
within a thousandth of each other, at \V{run.conv-small.score},
\V{run.transformer-small.score} and \V{run.transformer-medium.score}, and
each is behind the reference beyond the noise. Size does not change it: a
network of \V{run.transformer-small.parameters.millions} million parameters
and one of \V{run.transformer-medium.parameters.millions} million land in
the same place, and so does a body of another kind.
\ifvariants
Section~\ref{sec:changes} has the board-alone network that does pass the
trees: the same medium transformer, started from a body that was trained on
engine evaluations first.
\fi

The plain reading is that \V{recipe.games} games are too few for a network to
learn off the squares what the hand features count: a piece that can be
taken, the checks and captures on offer, the pressure on a king. Handed those
counts, the same medium transformer gains
\V{run.transformer-medium-hand.over-twin} over its twin that reads the board
alone, and passes the trees.
\ifpairs
Set against each other on the same resampled games, the one that reads the
counts is \V{pair.transformer-medium-hand.transformer-medium.verdict}
(Table~\ref{tab:pairs}).
\fi
So the body can use the facts. It did not find them by itself from this much
data.
\ifgames
What more games would buy is measured in Section~\ref{sec:changes}.
\else
More games is the run to add, since the recipe keeps a small part of one
month's.
\fi

Set against the ladder, the best run's gain is \V{best.gain.percent}\% of the
reference's own score, and \V{best.gain.of-lift}\% as large as the step from
the floor to the reference. Trees on \V{reference.features} numbers do most
of the work, and the network refines it.

\subsection{Every score}

\begin{table}[htbp]\centering\small
\block{metrics}
\caption{Every score of every run on the test positions. PR-AUC and ROC-AUC
are higher for a better run; the Brier score, the log loss and the
calibration error are lower, and are computed after calibration on the valid
rows.}
\label{tab:metrics}
\end{table}

Table~\ref{tab:metrics} has the other scores. ROC-AUC looks high for every
run that reads the evaluation, the trees with no board among them, and the
reason is not skill. A position that is already lost cannot be blundered:
with a win probability of \V{method.blunder-bar} or less, no move loses that
many more points. Every learner that reads the evaluation gets those rows
right for free, which lifts ROC-AUC for all of them alike. PR-AUC is spent on
the positions near the top of the ranking, where that gift does not help,
which is one more reason it is the score to read here.

\subsection{By the mover's rating}

Ranking is harder the stronger the mover. Every run's PR-AUC falls as the
band rises (Table~\ref{tab:bands}), and so does the floor's, which is the base
rate of each band: \V{band.1.floor} for movers rated \V{band.1.name} and
\V{band.5.floor} for those rated \V{band.5.name}. Stronger players blunder
less often, and a rarer event is harder to rank. So a band's column is read
from the top down, run against run, and never across.

\begin{table}[htbp]\centering\small
\block{bands-table}
\caption{PR-AUC by the mover's rating band. Read a column from the top down:
the base rate differs from band to band, and PR-AUC with it.}
\label{tab:bands}
\end{table}

Read that way, the best run leads the reference in \V{bands.best-ahead.of}.
Its lead is widest for movers rated \V{bands.lead.widest.name}, at
\V{bands.lead.widest}, and narrowest for those rated
\V{bands.lead.narrowest.name}, at \V{bands.lead.narrowest}
(Figure~\ref{fig:bands}). These are differences of scores without intervals,
and a band holds a fraction of the test games, so the order of the networks
inside one band is not to be leaned on.

\begin{figure}[htbp]\centering
\block{bands}
\caption{Each network's lead over the reference in PR-AUC, band by band: the
differences of the scores in Table~\ref{tab:bands}, without intervals. Above
the dashed line a network is ahead of the trees.}
\label{fig:bands}
\end{figure}

\iflook
The closer look makes the same cut with the comparison paired inside each
band (Table~\ref{tab:slice-band} in the appendix), and that is the one to
lean on. There \V{challenger.name} is ahead of the trees beyond the noise in
\V{slice.band.ahead.of}, the movers rated \V{slice.band.ahead.groups}, and
behind in \V{slice.band.behind}. Its gain is \V{slice.band.-1100.gain} for
movers under 1100 and \V{slice.band.2000+.gain} for those at 2000 and over,
where the interval holds zero. So what the network adds is proven for the
weaker players, who blunder most (\V{slice.band.-1100.happened} of their
moves, against \V{slice.band.2000+.happened} at 2000 and over) and whom a
coaching app is mostly for. For the strongest movers the test games do not
tell the two learners apart.
\fi

\subsection{A chance that can be read as one}

The app would show the number as a chance, so it has to mean what it says.
Table~\ref{tab:calibration} puts the test positions in
\V{best.calibration.groups} equal groups by a run's stated chance and sets
what the run said beside what happened. The groups run from
\V{best.calibration.lowest}\% to \V{best.calibration.highest}\%, and none is
further than \V{best.calibration.widest-gap} points from the rate that
happened (Figure~\ref{fig:calibration}). The correction fitted on the valid
rows has a slope of \V{best.platt.slope}. A slope of one is a network that
needed none, and one below it a network that stated its chances too far
apart.

\begin{table}[htbp]\centering\small
\block{calibration-table}
\caption{Each run's test positions in ten equal groups by its own stated
chance, lowest first: what the run said in the group on average, and the
share of the group that were blunders.}
\label{tab:calibration}
\end{table}

\begin{figure}[htbp]\centering
\block{calibration}
\caption{What was said against what happened, from the ten groups of
Table~\ref{tab:calibration}. On the diagonal a stated chance is the rate
that happened.}
\label{fig:calibration}
\end{figure}

The lowest groups are the positions that cannot be blundered: the stated
chance is nil there, and so is the rate. The highest group is what the model
can flag. Even there most moves are not blunders, which sets what a flag can
honestly say: this is a position where players like you go wrong unusually
often, never that you will.

\subsection{How the networks learned}

Figure~\ref{fig:curves} follows each network of the ladder on the valid
positions after every epoch. By that score, \V{ladder.at-last-epoch.of}
peaked at the last epoch given. The medium transformer with the hand
features had its best epoch at \V{run.transformer-medium-hand.best-epoch} of
\V{run.transformer-medium-hand.epochs}, so it had stopped improving before
its schedule ran out. Its twin on the board alone peaked at
\V{run.transformer-medium.best-epoch} of \V{run.transformer-medium.epochs}:
the network that has to find everything on the squares was still finding it.
\ifvariants
Section~\ref{sec:changes} gives both of them twice the epochs and says what
that bought.
\else
A run of each with twice the epochs says whether that matters, and it is the
cheapest run to add.
\fi

\begin{figure}[htbp]\centering
\block{curves}
\caption{Each ladder network's PR-AUC on the valid positions after every
epoch, before calibration. The dashed line is the reference's score on the
same rows.}
\label{fig:curves}
\end{figure}

\subsection{What the trees lean on}

The reference grew \V{reference.trees} trees over \V{reference.features}
features (Figure~\ref{fig:leans}). The three it leans on most are:
\V{reference.leans.1.feature} (\V{reference.leans.1.share}\% of the gain),
\V{reference.leans.2.feature} (\V{reference.leans.2.share}\%) and
\V{reference.leans.3.feature} (\V{reference.leans.3.share}\%). How the mover
stands comes first by a distance, which is the gift of the lost positions
again, and also a real pattern: Section~\ref{sec:closer} shows how sharply
the rate of blunders moves with it. A share of gain says what the trees used.
It does not say what makes a position dangerous.

\begin{figure}[htbp]\centering
\block{leans}
\caption{The ten features the reference's trees gained most from: each
feature's share of the total gain of the splits made on it.}
\label{fig:leans}
\end{figure}

\ifchanges
\section{One change at a time}
\label{sec:changes}

% One change at a time: the runs that change one thing in another run, set
% against that run, and how to read them. The same for every model, so a
% paper reads this under its own heading, after its results, then says what
% its own numbers mean, then reads studies.tex. Each part waits for its runs:
% a model that has none prints nothing.
\ifpairs
\paragraph{Runs set against the run each changed.}
The ladder says what each learner is worth. A change asks a narrower
question: what is one thing worth inside a learner? Twice the epochs, a body
that starts from another run's weights, the hand features under a network.
Each is a pair of runs that differ in that one thing.
Table~\ref{tab:pairs} sets the first of each pair against the second on the
same resampled test games, the way every run is set against the reference,
and gives the first run's gain. This is where the lab's rule for keeping a
change is read: a change is kept as a run of the model's own when its run is
better on the valid rows and the interval of its gain on the test rows
leaves zero out, however good its own score looks. A change that is not kept
is not thrown away: its row stays in this table, as a measurement of what
the change is worth.

\begin{table}[htbp]\centering\footnotesize
\block{pairs}
\caption{Runs set against the run each changed or has to beat. A gain is
the first run's over the second's, positive when the first is better. On the
valid rows it is the difference of the two scores. On the test rows it is
measured on the same resampled games, and it is bold when its 95\% interval
leaves zero out.}
\label{tab:pairs}
\end{table}
\fi

\ifvariants
Figure~\ref{fig:variants} follows each network of the ladder that another
run changed, beside the runs that changed it, on the valid rows. A run keeps
the colour of the network it changed, and its line is broken by what it
changed: long dashes for a longer schedule, dots for a start from another
run's body.

\begin{figure}[htbp]\centering
\block{variants}
\caption{A network of the ladder (the unbroken line) and the runs that
changed one thing in it (the broken lines), after every epoch on the valid
rows. One panel a network, all on one scale. The dashed level is the
reference's score on the same rows.}
\label{fig:variants}
\end{figure}
\fi


\ifvariants
\paragraph{What twice the epochs were worth here.}
Each medium transformer was trained again on a schedule of
\V{run.transformer-medium-long.epochs} epochs for the ladder's
\V{run.transformer-medium.epochs}. On the board alone,
\V{run.transformer-medium-long.name} is
\V{pair.transformer-medium-long.transformer-medium.valid.verdict} than the
run it changed on the valid positions and
\V{pair.transformer-medium-long.transformer-medium.verdict} on the test
positions (\V{pair.transformer-medium-long.transformer-medium.gain},
\V{pair.transformer-medium-long.transformer-medium.gain.low} to
\V{pair.transformer-medium-long.transformer-medium.gain.high}): it is
\V{pair.transformer-medium-long.transformer-medium.kept}. With the hand
features, \V{run.transformer-medium-hand-long.name} is
\V{pair.transformer-medium-hand-long.transformer-medium-hand.valid.verdict}
than the run it changed on the valid positions and
\V{pair.transformer-medium-hand-long.transformer-medium-hand.verdict} on the
test positions
(\V{pair.transformer-medium-hand-long.transformer-medium-hand.gain},
\V{pair.transformer-medium-hand-long.transformer-medium-hand.gain.low} to
\V{pair.transformer-medium-hand-long.transformer-medium-hand.gain.high}): it
is \V{pair.transformer-medium-hand-long.transformer-medium-hand.kept}
either. So neither longer run takes the place of the run it changed. A run
that is not kept leaves the leaderboard when it is filed, and its row stays
in Table~\ref{tab:pairs} as the measurement of what the epochs were worth.

Neither run used its whole schedule. Both had their best valid epoch at
\V{run.transformer-medium-long.best-epoch} of
\V{run.transformer-medium-long.epochs}, and each was stopped
\V{networks.patience} epochs later, as the trainer stops any run whose valid
score has not improved for that long (Figure~\ref{fig:variants}). The valid
positions come from few games, and their score moves from one epoch to the
next by more than the two schedules differ, so where a run stopped is partly
chance. As trained, the longer runs are no better.

A reading: this model is short of games, not of steps. The transformer on
the board alone was still improving when the ladder's schedule ran out, and
more of the same games did not carry it further. That fits the plain reading
of Section~\ref{sec:results}, that \V{recipe.games} games are too few for
a network to learn off the squares what the hand features count.

\paragraph{What a start from engine evaluations was worth here.}
The lab's position-evaluation model trains the same medium transformer on
positions from the open Lichess evaluations file, to say the mover's winning
chances as an engine sees them. That label says nothing about people. The
model is there for its body: a run of this model can start from it. The
embeddings and the blocks are copied in, the head starts new, and the run
then trains on this model's games on the ladder's schedule, with nothing
else changed.

On the board alone, \V{run.transformer-medium-pretrained.name} is
\V{pair.transformer-medium-pretrained.transformer-medium.valid.verdict} than
the run it changed on the valid positions and
\V{pair.transformer-medium-pretrained.transformer-medium.verdict} on the
test positions (\V{pair.transformer-medium-pretrained.transformer-medium.gain},
\V{pair.transformer-medium-pretrained.transformer-medium.gain.low} to
\V{pair.transformer-medium-pretrained.transformer-medium.gain.high}): it is
\V{pair.transformer-medium-pretrained.transformer-medium.kept}. With the
hand features, \V{run.transformer-medium-hand-pretrained.name} is
\V{pair.transformer-medium-hand-pretrained.transformer-medium-hand.valid.verdict}
than the run it changed on the valid positions and
\V{pair.transformer-medium-hand-pretrained.transformer-medium-hand.verdict}
on the test positions
(\V{pair.transformer-medium-hand-pretrained.transformer-medium-hand.gain},
\V{pair.transformer-medium-hand-pretrained.transformer-medium-hand.gain.low}
to
\V{pair.transformer-medium-hand-pretrained.transformer-medium-hand.gain.high}):
it is
\V{pair.transformer-medium-hand-pretrained.transformer-medium-hand.kept}.
The model's best run is \V{best.name}.

Two things follow. The first is that the start did what neither size nor
steps had done. From the board alone the transformer now passes the trees,
with a gain over the reference of \V{run.transformer-medium-pretrained.gain}
(\V{run.transformer-medium-pretrained.gain.low} to
\V{run.transformer-medium-pretrained.gain.high}), where the same network
started from nothing has one of \V{run.transformer-medium.gain}. The second
is that the counts add little once the body has been trained first.
\V{run.transformer-medium-hand-pretrained.name} scores
\V{run.transformer-medium-hand-pretrained.over-twin} above
\V{run.transformer-medium-pretrained.name}, where on the ladder the same
difference was \V{run.transformer-medium-hand.over-twin}, and the two runs'
gains over the reference lie inside each other's intervals
(Table~\ref{tab:leaderboard}). The two were not set against each other on
resampled games, so that difference is a reading of the leaderboard.

A reading: this is the shortage of games again, met from another side. A
body that has read a great many positions for another label arrives able to
read a board, and most of what the hand features count comes with it.
Neither run was at its best in its last epoch: the best valid epochs were
\V{run.transformer-medium-pretrained.best-epoch} and
\V{run.transformer-medium-hand-pretrained.best-epoch} of
\V{run.transformer-medium-hand-pretrained.epochs}, so a longer schedule from
the same start is not the next run.
\fi

\ifshapes
\paragraph{What another shape of network was worth here.}
The architecture test keeps a medium transformer's rows, its head and its
epochs, and changes its shape. Each run below started from nothing, as the
network it changes did. Each is set against that network in
Table~\ref{tab:pairs} and followed epoch by epoch in
Figure~\ref{fig:shapes}.
\begin{itemize}
\item \emph{A convolutional tower}, \V{run.conv-medium-hand.name}: residual
blocks of convolutions, each also weighing its channels by a mean over the
whole board, with \V{run.conv-medium-hand.parameters.millions} million
parameters for the transformer's
\V{run.transformer-medium-hand.parameters.millions} million. It is trained at a
tower's peak learning rate, \V{run.conv-medium-hand.lr} for
\V{run.transformer-medium-hand.lr}, and changes nothing else. It is
\V{pair.conv-medium-hand.transformer-medium-hand.valid.verdict} than
\V{run.transformer-medium-hand.name} on the valid positions and
\V{pair.conv-medium-hand.transformer-medium-hand.verdict} on the test
positions (\V{pair.conv-medium-hand.transformer-medium-hand.gain},
\V{pair.conv-medium-hand.transformer-medium-hand.gain.low} to
\V{pair.conv-medium-hand.transformer-medium-hand.gain.high}): it is
\V{pair.conv-medium-hand.transformer-medium-hand.kept}.
\item \emph{A hybrid}, \V{run.hybrid-medium-hand.name}: a few residual blocks
of convolutions first, then layers of the transformer over the squares
they give, \V{run.hybrid-medium-hand.parameters.millions} million parameters
in all. It is
\V{pair.hybrid-medium-hand.transformer-medium-hand.valid.verdict} than
\V{run.transformer-medium-hand.name} on the valid positions and
\V{pair.hybrid-medium-hand.transformer-medium-hand.verdict} on the test
positions (\V{pair.hybrid-medium-hand.transformer-medium-hand.gain},
\V{pair.hybrid-medium-hand.transformer-medium-hand.gain.low} to
\V{pair.hybrid-medium-hand.transformer-medium-hand.gain.high}): it is
\V{pair.hybrid-medium-hand.transformer-medium-hand.kept}.
\item \emph{Attention told what each piece attacks.} The offsets say that
two squares share a diagonal, and not whether anything stands between
them. Here each head also learns a bias for a pair of squares when the
piece on one attacks the other on this board, worked out from the tokens
with no search. With the hand features, \V{run.transformer-medium-hand-attacks.name} is
\V{pair.transformer-medium-hand-attacks.transformer-medium-hand.valid.verdict} than
\V{run.transformer-medium-hand.name} on the valid positions and
\V{pair.transformer-medium-hand-attacks.transformer-medium-hand.verdict} on the test
positions (\V{pair.transformer-medium-hand-attacks.transformer-medium-hand.gain},
\V{pair.transformer-medium-hand-attacks.transformer-medium-hand.gain.low} to
\V{pair.transformer-medium-hand-attacks.transformer-medium-hand.gain.high}): it is
\V{pair.transformer-medium-hand-attacks.transformer-medium-hand.kept}. On the
board alone, \V{run.transformer-medium-attacks.name} is
\V{pair.transformer-medium-attacks.transformer-medium.valid.verdict} than
\V{run.transformer-medium.name} on the valid positions and
\V{pair.transformer-medium-attacks.transformer-medium.verdict} on the test
positions (\V{pair.transformer-medium-attacks.transformer-medium.gain},
\V{pair.transformer-medium-attacks.transformer-medium.gain.low} to
\V{pair.transformer-medium-attacks.transformer-medium.gain.high}): it is
\V{pair.transformer-medium-attacks.transformer-medium.kept}.
\item \emph{Attention without the offsets},
\V{run.transformer-medium-hand-no-geometry.name}: a study, there to measure what
the offsets are worth and never a candidate. Without them the transformer
is \V{pair.transformer-medium-hand-no-geometry.transformer-medium-hand.valid.verdict}
on the valid positions than with them, and
\V{pair.transformer-medium-hand-no-geometry.transformer-medium-hand.verdict} on the
test positions (\V{pair.transformer-medium-hand-no-geometry.transformer-medium-hand.gain},
\V{pair.transformer-medium-hand-no-geometry.transformer-medium-hand.gain.low} to
\V{pair.transformer-medium-hand-no-geometry.transformer-medium-hand.gain.high}).
\end{itemize}

No run of another shape is kept. On the test positions every one of them is
inside the noise of the transformer it changes, so the model's runs stand as
they did and its best run is still \V{best.name}. The tower, and the
attention that is told the attacks on the board alone, pass the first half
of the lab's rule, on the valid positions, and not the second.

A reading: what this model's transformer lacks is not a way of reading the
board. Told what each piece attacks, the transformer on the board alone
still trails the trees, with a gain over the reference of
\V{run.transformer-medium-attacks.gain}
(\V{run.transformer-medium-attacks.gain.low} to
\V{run.transformer-medium-attacks.gain.high}). The lab's own paper sets
this beside puzzle difficulty, where the same change is kept.
\iffeatures
The features study below says where the difference may lie. What the counts
give this network is mostly the material group, three sums over the whole
board, and which piece attacks which is not what those turn on.
\fi
It fits the reading of the longer schedule as well: this model is short of
games, and of the changes made to its transformer the one that moved it is
the start from a body that had read many more positions elsewhere.
\fi

% The two studies every model has: fewer training games, and one group of
% hand features left out at a time. The same for every model, so a paper
% reads this after changes.tex and what it says of its own runs there, and
% adds what its own numbers mean. Each study waits for its runs.
\ifgames
\paragraph{Fewer training games.}
How much of a score is the data, and would more of it help? The games study
trains a run again on half and on a quarter of the training games, and
scores it on the same test rows (Table~\ref{tab:games},
Figure~\ref{fig:games}). The games are chosen by a hash, so the quarter lies
inside the half. A network is given twice the epochs on half the games and
four times on a quarter, so it takes as many steps each time and loses only
the games. The study is read forward: where halving the games costs little,
doubling them would buy little, so the curve says what a larger dataset is
worth before anyone builds one.

\begin{table}[htbp]\centering\footnotesize
\block{study-games}
\caption{The games study. Each run is trained on a share of the training
games and scored on all the test rows. Its gain is paired with the run that
was trained on all the games: below zero the run is worse, and bold when the
95\% interval leaves zero out.}
\label{tab:games}
\end{table}

\begin{figure}[htbp]\centering
\block{study-games-figure}
\caption{The games study as a learning curve: the score on the test rows by
the share of the training games a run was given. Each step to the right
doubles the games.}
\label{fig:games}
\end{figure}
\fi

\iffeatures
\paragraph{A group of hand features left out.}
The hand features come in groups: material, the moves on offer, pieces that
can be taken, passed pawns and attacks near the kings, and for a model that
reads a line of moves, its first move and the line as a whole. The features
study trains a run again without one group at a time
(Table~\ref{tab:features}, Figure~\ref{fig:features}). A run that is worse
without a group was using it, and nothing else the run reads holds the same.
A run that is no worse either made no use of the group or finds the same
elsewhere, in another group or on the board. So the study measures what a
group adds on top of everything else, which is less than what it knows: a
group can matter a great deal on a board and cost nothing here.

\begin{table}[htbp]\centering\footnotesize
\block{study-features}
\caption{The features study. Each run reads every hand feature but one
group. Its gain is paired with the run that reads them all: below zero the
run is worse without the group, and bold when the 95\% interval leaves zero
out. The group whose loss costs most comes first.}
\label{tab:features}
\end{table}

\begin{figure}[htbp]\centering
\block{study-features-figure}
\caption{What leaving a group of hand features out costs the reference: the
paired gain of the run without the group over the run with it, with its 95\%
interval. Left of the line the run is worse without the group.}
\label{fig:features}
\end{figure}
\fi


\ifgames
\paragraph{What more games are worth here.}
The network is worse on fewer games in
\V{study.games.transformer-medium-hand.worse.of}. Against the run trained on every training
game, \V{run.transformer-medium-hand.name} has a paired gain of \V{pair.transformer-medium-hand-half.transformer-medium-hand.gain} on
half of them and of \V{pair.transformer-medium-hand-quarter.transformer-medium-hand.gain} on a quarter. The trees are
worse beyond the noise in \V{study.games.gbm-hand.worse.of}, with gains
of \V{pair.gbm-hand-half.gbm-hand.gain} and
\V{pair.gbm-hand-quarter.gbm-hand.gain}. On a quarter of the games the two
learners are level, and the network's lead opens as the games are added
(Table~\ref{tab:games}). That is the reading of the longer schedule,
now measured: this model is short of games. \V{recipe.games} games of one
month are a small part of what the Lichess database holds, so a larger
dataset is the next thing to try.
\fi

\ifmore
\paragraph{What more games were worth.}
That reading was then checked on a larger build of the dataset, which holds
\V{more.gbm-hand.most} the training games and every row these runs are
scored on. Against the run trained on today's games,
\V{run.transformer-medium-hand.name} has a paired gain of
\V{pair.transformer-medium-hand-double.transformer-medium-hand.gain} on twice the games and of
\V{pair.transformer-medium-hand-quadruple.transformer-medium-hand.gain} on four times, which is
\V{pair.transformer-medium-hand-quadruple.transformer-medium-hand.verdict}. The trees' gains are
\V{pair.gbm-hand-double.gbm-hand.gain} and
\V{pair.gbm-hand-quadruple.gbm-hand.gain}, the second
\V{pair.gbm-hand-quadruple.gbm-hand.verdict}. By the lab's rule the
network on four times the games is \V{more.transformer-medium-hand.most.kept} and the trees
on four times the games are \V{more.gbm-hand.most.kept}.
\fi

\iffeatures
\paragraph{Which counts each learner needs here.}
The study was run twice here: on the trees, and on the transformer that
reads the hand features. Both lose most without \V{study.features.gbm-hand.most.group}: the trees' paired gain
without it is \V{study.features.gbm-hand.most.gain} and the network's
\V{study.features.transformer-medium-hand.most.gain}. For the network that is most of what
the hand features gave it in the first place
(\V{pair.transformer-medium-hand.transformer-medium.gain}, Table~\ref{tab:pairs}). Loose pieces
and the moves on offer cost both learners something: the trees
\V{pair.gbm-hand-no-hanging.gbm-hand.gain} and
\V{pair.gbm-hand-no-mobility.gbm-hand.gain}, the network
\V{pair.transformer-medium-hand-no-hanging.transformer-medium-hand.gain} and \V{pair.transformer-medium-hand-no-mobility.transformer-medium-hand.gain}. The
attackers near the kings split the two. Without them the trees are
\V{pair.gbm-hand-no-kings.gbm-hand.verdict}
(\V{pair.gbm-hand-no-kings.gbm-hand.gain}) and the network is
\V{pair.transformer-medium-hand-no-kings.transformer-medium-hand.verdict}. Without the passed pawns the trees
are \V{pair.gbm-hand-no-pawns.gbm-hand.verdict} and the network is
\V{pair.transformer-medium-hand-no-pawns.transformer-medium-hand.verdict}.

A reading: the network reads off the squares what stands near a king and
where the pawns are. What it does not learn from \V{recipe.games} games is
the three sums of the material group: how much is on the board, who is
ahead, and how far the game has gone. Those stay worth handing to it.
\fi
\fi

\section{A closer look}
\label{sec:closer}

% What a closer look is and how it is made: the same for every model. A paper
% reads it under its own section heading, before the cuts that are its own.
\iflook
A leaderboard says which run is best. It does not say where a model is right
and where it is wrong, and one score is a poor guide to whether the app can
use it. So the same test rows are cut several ways, on the machine that has
the dataset and every run's test predictions, and the small tables that come
out are committed beside the leaderboard. Everything in this section and the
next is drawn from them (\V{look.test-rows} test rows, looked at on
\V{look.date}).

\paragraph{Slices.}
The test rows are grouped by one column at a time. For each group a slice
gives how many rows fall in it, what happened there, what a run said on
average, and how the run scored. Two runs are set side by side: \V{challenger.name},
the best run that is not the reference, and the reference. The gain of the
first over the second is measured inside each group as it is for the whole
test set, paired and over resampled games (\V{look.resamples} draws), and is
left out for a group with fewer than \V{look.smallest-compared} test rows.
A slice is where a claim such as ``the network helps most on long lines'' is
checked, not made: a group whose interval holds zero says nothing either way.

\paragraph{Grids.}
Two columns at once: what happened in each cell, and what the run said there,
on one scale of colour so the two can be compared by eye. A cell with fewer
than \V{grid.smallest-cell} test rows is left empty.

\paragraph{Examples.}
Positions are chosen by rule, never by eye: ordinary rows picked by a hash of
their key, the rows the best run is surest of at each end, and its largest
misses each way, \V{examples.count} of each. They show what the numbers are
about. They are not evidence of how often anything happens.
\else
\emph{This build of the paper was drawn before the model's closer look was
run. The closer look cuts the test rows by rating, clock and the shape of a
position, checks what a stated number is worth, and shows example positions.
It is written by \texttt{lab analyse} on the machine that holds the dataset
and every run's test predictions, and this section and the next appear once
its file is committed.}
\fi


\iflook
\subsection{What a flag is worth}

The app would not show every chance. It would flag the positions at the top
of the ranking, so the question is what the top holds.
Table~\ref{tab:ranking-top} takes a share of the test positions from the
highest stated chances down. Of the hundredth the best run calls most
dangerous, \V{best.top.1.rate}\% were blunders, \V{best.top.1.lift} times the
base rate, and that hundredth holds \V{best.top.1.recall}\% of all the
blunders in the test games. Of the tenth it calls most dangerous,
\V{best.top.10.rate}\% were blunders, and that tenth holds
\V{best.top.10.recall}\% of them.

\begin{table}[htbp]\centering\small
\block{ranking-top}
\caption{The top of each run's ranking. For a share of the test positions
taken from the highest stated chances: the share of them that were blunders
(\emph{rate}), the share of all blunders that fall among them (\emph{of
all}), and the rate against the base rate (\emph{lift}).}
\label{tab:ranking-top}
\end{table}

The two ends of the table are two ways to use a flag, and the app has to
choose between them. A rare flag is right often and passes over most
blunders. A frequent one catches most of them and is wrong most of the time:
of the fifth of the positions called most dangerous, \V{best.top.20.rate}\%
were blunders, and that fifth holds \V{best.top.20.recall}\% of them. No
setting does both. In a review a player reads a few flags a game, which is
the rare end, and the table says what those few are worth.

Table~\ref{tab:ranking-bottom} reads the ranking from the other end: of the
positions a run calls safest, how many were indeed not blunders. Of the half
the best run calls safest, \V{best.bottom.50.rate}\% were not, and the other
half of the ranking holds \V{best.top.50.recall}\% of all the blunders. Part
of that is given away for nothing: the positions that are already lost,
which no move can blunder, are about a fifth of the test positions (the two
lowest groups of Table~\ref{tab:win}) and sit at the bottom of every
ranking.

\begin{table}[htbp]\centering\small
\block{ranking-bottom}
\caption{The bottom of each run's ranking. For a share of the test positions
taken from the lowest stated chances: the share of them that were not
blunders, the share of all moves that were not blunders that fall among them,
and that rate against the rate over every test position.}
\label{tab:ranking-bottom}
\end{table}

\subsection{By how the mover stands}

Table~\ref{tab:win} cuts the test positions by the mover's winning chances
before the move, and Figure~\ref{fig:win} sets the rate of blunders in each
group beside what the two runs said. Below \V{method.blunder-bar}\% no move
can lose \V{method.blunder-bar} points more, so the two lowest groups,
\V{slice.win.-10.share}\% and \V{slice.win.10-20.share}\% of the test
positions, hold no blunders, and both runs say so.

Above that, the rate is not highest where the game is in the balance. It is
highest where the mover is winning and has not yet won:
\V{slice.win.80-90.happened} at 80 to 90\% and \V{slice.win.60-80.happened}
at 60 to 80\%, against \V{slice.win.40-60.happened} at 40 to 60\% and
\V{slice.win.20-40.happened} at 20 to 40\%. Past 90\% it falls again, to
\V{slice.win.90+.happened}. An advantage is the dangerous thing to hold.

One reading of that is the win not taken. A position can stand at 85\%
because the opponent has just gone wrong, and then the advantage is there
only for a mover who finds the move that punishes it. Playing something else
hands the points back, and the label counts that as a blunder like any
other. The tables cannot tell a piece left to be taken from a win not taken,
since the label is one for both, and how much of the peak is each is a cut
nobody has made.

\begin{table}[htbp]\centering\footnotesize
\block{slice-win}
\caption{The test positions by the mover's winning chances before the move.
For each group: its test rows, the share of moves that were blunders, what
each of the two runs said on average and its PR-AUC there, and the first
run's paired gain over the second.}
\label{tab:win}
\end{table}

\begin{figure}[htbp]\centering
\block{slice-win-said}
\caption{The share of moves that were blunders by the mover's winning chances
before the move, beside what the two runs said on average.}
\label{fig:win}
\end{figure}

Both runs follow the rate from group to group: at its furthest
\V{challenger.name} is \V{slice.win.challenger.widest-off} points from what
happened, and the reference \V{slice.win.reference.widest-off}. Among the
groups with enough rows to compare, the network is ahead of the trees beyond
the noise in \V{slice.win.ahead.of} (\V{slice.win.ahead.groups}) and behind
in \V{slice.win.behind}.

\subsection{By the clock}

The clock is the part of the question an engine has no view of.
Table~\ref{tab:clock} cuts the test positions by the seconds on the mover's
clock, and Figure~\ref{fig:clock} shows the rate of blunders in each group
beside what the runs said.

\begin{table}[htbp]\centering\footnotesize
\block{slice-clock}
\caption{The test positions by the seconds on the mover's clock.}
\label{tab:clock}
\end{table}

\begin{figure}[htbp]\centering
\block{slice-clock-said}
\caption{The share of moves that were blunders by the seconds on the mover's
clock, beside what the two runs said on average.}
\label{fig:clock}
\end{figure}

The first thing the table says is how little of a rapid game is played short
of time. \V{slice.clock.480+.share}\% of the test positions have eight
minutes or more on the mover's clock, and the two groups under a minute hold
\V{slice.clock.-30.rows} and \V{slice.clock.30-60.rows} of the
\V{look.test-rows}. In those two groups and the next, the rate of blunders
goes \V{slice.clock.-30.happened}, \V{slice.clock.30-60.happened},
\V{slice.clock.60-120.happened}: with so few positions that is noise as much
as pattern, and no difference between the runs there is beyond it. Where the
positions are, the network's lead is proven: it is ahead of the trees beyond
the noise in \V{slice.clock.ahead.of} (\V{slice.clock.ahead.groups}).

So this dataset cannot say what a short clock does, which is the thing a
player most wants to know. The cut is also confounded. A short clock comes
late in a game, and late positions are more often already decided, where
little is left to lose (Table~\ref{tab:slice-ply}). A learner reads the
clock, the ply and the evaluation together and can hold them apart. A table
by the clock alone cannot. Time trouble needs a dataset with more of it:
many more rapid games, or blitz.

Table~\ref{tab:seconds} uses something no learner reads, because it is known
only once the move is made: how long the move took.

\begin{table}[htbp]\centering\footnotesize
\block{slice-seconds}
\caption{The test positions by the seconds the move took, which no learner
reads.}
\label{tab:seconds}
\end{table}

The rate of blunders rises with every step: \V{slice.seconds.-2.happened} of
the moves made in under two seconds, \V{slice.seconds.5-10.happened} of
those that took five to ten, and \V{slice.seconds.40+.happened} of those
that took forty or more. A long think is not a safe think. It marks a
position the player found hard, and hard positions are where blunders are
made.

Neither run knows how long the move will take, and both are off in the same
direction at both ends. For moves made in under two seconds
\V{challenger.name} states a chance \V{slice.seconds.-2.challenger.off}
points from the rate that happened, and for the longest thinks
\V{slice.seconds.40+.challenger.off} points, where a figure below zero is a
chance stated too low. The reference is off by
\V{slice.seconds.-2.reference.off} and \V{slice.seconds.40+.reference.off}.
So the danger the models understate lies in the slow moves, not the fast
ones, and the time a player takes knows something about the position that
neither learner reads off it. That time is the label of another model of
this lab, which predicts the seconds from the same inputs. Whether one body
trained on both labels would close the gap is a run to make.

\subsection{By what can be taken, and how much there is to choose from}

A piece that can be taken is the plainest danger on a board, and the hand
features count it for the trees. Table~\ref{tab:hanging} cuts the test
positions by how many of the mover's own pieces can be taken.

\begin{table}[htbp]\centering\footnotesize
\block{slice-hanging}
\caption{The test positions by how many of the mover's pieces can be taken:
pieces that something cheaper attacks, or that anything attacks while nothing
defends them.}
\label{tab:hanging}
\end{table}

The rate of blunders rises with the count, and by less than a player might
expect: from \V{slice.hanging.0.happened} with no such piece to
\V{slice.hanging.3+.happened} with three or more. A piece under attack is a
danger a player sees. \V{slice.hanging.0.share}\% of the test positions have
none, and more than half of all the blunders are made there. The opponent's
pieces that can be taken tell even less, from
\V{slice.opponent-hanging.0.happened} with none to
\V{slice.opponent-hanging.3+.happened} with three or more
(Table~\ref{tab:slice-opponent-hanging}).

What moves the rate more is how much there is to choose from. A mover in
check blunders \V{slice.check.yes.happened} of the time, against
\V{slice.check.no.happened} for one who is not
(Table~\ref{tab:slice-check}): a check leaves few legal moves. By the count
of legal moves the rate climbs from \V{slice.moves.-5.happened} with fewer
than five to \V{slice.moves.45+.happened} with 45 or more
(Table~\ref{tab:slice-moves}), and it is in the groups with the most moves
that the network leads the trees beyond the noise
(\V{slice.moves.ahead.groups}). Many legal moves go with a full board and a
middlegame, so this cut does not hold choice apart from the stage of the
game either.

\subsection{By the stage of the game}

Two cuts in the appendix follow a game from its opening to its ending: by
the ply (Table~\ref{tab:slice-ply}) and by the knights, bishops, rooks and
queens left on the board (Table~\ref{tab:slice-pieces}). The rate of
blunders is highest between ply 30 and 40, at \V{slice.ply.30-40.happened},
and lowest from ply 80 on, at \V{slice.ply.80+.happened}.

The network's lead is proven early in a game and not late.
\V{challenger.name} is ahead beyond the noise in \V{slice.ply.ahead.of}
(\V{slice.ply.ahead.groups}) and behind beyond the noise in
\V{slice.ply.behind.of}. From ply 60 to 80 its gain is
\V{slice.ply.60-80.gain} and from ply 80 on \V{slice.ply.80+.gain}, where a
gain below zero puts the trees ahead, and both intervals hold zero. The cut
by material agrees as far as it can: with fewer than nine pieces left the
network's gain is inside the noise in all three groups
(\V{slice.pieces.-3.gain}, \V{slice.pieces.3-6.gain} and
\V{slice.pieces.6-9.gain}), and it is ahead beyond the noise with nine or
more (\V{slice.pieces.ahead.groups}).

So the network earns its lead in the opening and the middlegame, and in
long endings the test games do not tell the two learners apart. Positions
from ply 80 on are \V{slice.ply.80+.share}\% of the test positions and as
small a share of the ones a network trains on, so a difference there would
have to be large to show. For the app that leaves the trees as the safe
choice late in a game: nothing here proves the network better there.

\subsection{Two things at once}

Figure~\ref{fig:grid-win} crosses the mover's rating with how the mover
stands, and Figure~\ref{fig:grid-clock} crosses it with the clock. In each,
the upper grid is the rate of blunders that happened and the lower one what
the run said. If the lower grid looks like the upper one, the run has the
joint pattern and not only each margin.

\begin{figure}[htbp]\centering
\block{grid-rating-win}

\medskip
\block{grid-rating-win-said}
\caption{The share of moves that were blunders, in percent, by the mover's
rating and the mover's winning chances before the move. Above: what happened.
Below: what the best run that is not the reference said. Both grids share one
scale of colour, and a cell with too few test positions is left empty.}
\label{fig:grid-win}
\end{figure}

The first grid is the clearest picture in this paper of what the model is
for. Movers under 1100 who stand at 60 to 80\% blunder
\V{grid.rating-win.-1100.60-80.happened} of the time, the highest cell of
the grid. Movers at 2000 and over in the same spot blunder
\V{grid.rating-win.2000+.60-80.happened} of the time. The run said
\V{grid.rating-win.-1100.60-80.said} and
\V{grid.rating-win.2000+.60-80.said}. An advantage is several times as
dangerous in a weaker player's hands, which neither margin states by itself.
The model has the pattern: it puts the first far above the second, as the
truth does.

\begin{figure}[htbp]\centering
\block{grid-rating-clock}

\medskip
\block{grid-rating-clock-said}
\caption{The share of moves that were blunders, in percent, by the mover's
rating and the seconds on the mover's clock. Above: what happened. Below:
what the run said.}
\label{fig:grid-clock}
\end{figure}

The clock grid is rougher, for the reason the clock's own table gave. Its
short-clock columns hold few positions a cell, and what happened there jumps
about: for movers under 1100, \V{grid.rating-clock.-1100.30-60.happened}
with 30 to 60 seconds on the clock (\V{grid.rating-clock.-1100.30-60.rows}
positions) and \V{grid.rating-clock.-1100.60-120.happened} with 60 to 120
(\V{grid.rating-clock.-1100.60-120.rows}). What the run says stays smooth,
\V{grid.rating-clock.-1100.30-60.said} and
\V{grid.rating-clock.-1100.60-120.said}. With so few positions in a cell
neither grid proves the other wrong.

\subsection{Beside the reference}

Table~\ref{tab:beside} sets each run's stated chances beside the reference's
on the same test positions. The best run's correlation with the reference is
\V{best.beside.correlation}, and it is the closer of the two to what happened
on \V{best.beside.closer}\% of the positions. Two runs that score alike and
still disagree position by position would do better averaged than either
does alone, and the stage of the game says where the network's lead is
proven. A blend of the two is a run nobody has made yet.

\begin{table}[htbp]\centering\small
\block{beside}
\caption{Each run's stated chances beside the reference's on the same test
positions: their correlation, and how often the run is the closer of the two
to what happened.}
\label{tab:beside}
\end{table}

\section{Positions}
\label{sec:positions}

The positions below were chosen by rule from the test games. Each board is
seen from the side to move. The grey arrow is the opponent's last move and
the blue arrow is the move that was played. Beside each board are the move,
whether it was a blunder, what the best run and the reference said before it,
and what the position came with. The games are public Lichess games; no game
is named.

The surprises are the instructive ones. Figure~\ref{fig:surprised-no} holds
the positions the model called most dangerous in which the mover did not
blunder, and Figure~\ref{fig:surprised-yes} the blunders it saw coming least.
The positions it is surest of at each end are in the appendix
(Figures~\ref{fig:highest} and~\ref{fig:lowest}).

\examples{examples-typical}{Ordinary positions: the first \V{examples.count}
test positions in the order of a hash of their key, so nothing about them
picked them.}{fig:typical}

\examples{examples-surprised-no}{Called dangerous, and the mover did not
blunder: the test positions with the highest stated chance among the moves
that were not blunders.}{fig:surprised-no}

\examples{examples-surprised-yes}{Called safe, and the mover blundered: the
test positions with the lowest stated chance among the moves that were
blunders.}{fig:surprised-yes}
\fi

\section{What the app can use today}
\label{sec:use}

\begin{itemize}
\item \textbf{A chance that means what it says.} After calibration none of
the best run's groups is further than \V{best.calibration.widest-gap} points
from the rate that happened, so on Lichess rapid games the number can be
shown as a chance.
\iflook
\item \textbf{A flag for the positions worth slowing down in.} Flagged at
the top hundredth of the ranking, a position is one where
\V{best.top.1.rate}\% of the moves played were blunders, against
\V{label.rate.test}\% over all the test positions. The words that go with
such a flag are that players of this level often go wrong here, never that
this player will.
\item \textbf{A lapse or a hard position.} After a blunder, the chance
stated before it says which it was. At a chance near the base rate the
player slipped where few do. At several times the base rate the position
trips most players of that level, and the review can say so.
\item \textbf{Most for the weaker player.} The network's lead over the trees
is proven for the movers rated \V{slice.band.ahead.groups}. The two weakest
bands are among them, and those are the players who blunder most.
\item \textbf{The trees first, and late in a game.} \V{task.reference}
reaches \V{reference.score} against the best run's \V{best.score}. It is a
file of trees that needs no GPU, and from ply 60 on the test games do not
prove the network better. Where the app has to run the model itself, the
trees are the place to start.
\item \textbf{Not yet a model of time trouble.} Too few of these positions
have a short clock for the tables to say what one does.
\else
\item \textbf{The trees first.} \V{task.reference} reaches
\V{reference.score} against the best run's \V{best.score}. It is a file of
trees that needs no GPU, so where the app has to run the model itself, the
trees are the place to start.
\fi
\end{itemize}

\section{What the model does not say yet}
\label{sec:limits}

\begin{itemize}
\item \textbf{Lost positions are free points.} A position that is already
lost cannot be blundered, and every learner that reads the evaluation knows
it. That lifts every run's ROC-AUC alike. A run on the positions that can
still be blundered would be the fairer picture of each learner, and is its
own task.
\item \textbf{A blunder is two things under one label.} A piece left to be
taken and a win not taken both lose the points. They are different mistakes
with different lessons, and telling them apart is a label of its own.
\item \textbf{Only analysed games give a label.} Lichess analyses a game when
a player asks for it, and players choose which games to analyse.
\item \textbf{The label is as deep as the server's analysis.} A move the
server's engine misjudged is mislabelled here too.
\item \textbf{One position, no history.} The model does not see the moves or
the clock use that led to the position, which a model of the game as a
sequence would.
\item \textbf{Short clocks are scarce.} A rapid game is mostly played with
time in hand, so the test games hold too few positions under a minute to
say what a short clock does.
\item \textbf{Lichess ratings, rapid games, one month.} The chance is on the
scale of Lichess rapid ratings. Stated for a player on another site, it needs
calibrating to that site's ratings first, and a check on the player's own
games.
\item \textbf{A piece that can be taken is a proxy.} The hand feature counts
attackers and defenders. It is not a full search of the exchanges on the
square.
\end{itemize}

The lab's first hazard model was a tree model over sixteen board features,
trained on rated games of April 2019 with scripts of its own. It showed that
the size of the evaluation alone ranks blunders no better than chance and
that a few board features do. Its data, its features and its split differ
from this harness's, so its scores are not comparable with the ones here.

\section{Cost and how to draw this again}
\label{sec:cost}

The \V{runs.count} runs took \V{cost.seconds} in all: \V{cost.seconds.networks}
for the networks, which trained on a rented GPU (\V{best.device}), and
\V{cost.seconds.others} for the floor and the trees. The best run has
\V{best.parameters} parameters and trained for \V{best.epochs.run} epochs in
\V{best.seconds}, at about \V{best.rows-a-second} positions a second.
Table~\ref{tab:runs} has every run's time and Table~\ref{tab:settings} the
ladder networks' settings.

Everything is drawn from committed files. The dataset is built by
\texttt{lab data game-positions} from the recipe, the runs are trained by
\texttt{lab all blunder-hazard}, each run writes its record under the model's
\texttt{results/} folder, and \texttt{lab paper blunder-hazard} writes every
number, table and figure of this paper from those records and compiles it. A
record names the build of the data it read (\V{data.fingerprint}) and the
commit of the code that trained it (\V{code.commit}). Weights and test
predictions are not committed. They are kept on the owner's machine, and the
model's card says where and gives the hash of the best run's weights.

\section{What comes next}

\begin{itemize}
\ifgames\else
\item \textbf{More games.} No network learned from the squares what the hand
features count. The recipe keeps a small part of one month's games, so a
larger build is the first run to add, and a run on half and on a quarter of
these games says beforehand what it would buy.
\fi
\iffeatures\else
\item \textbf{Which hand features carry it.} Leaving out one group of
features at a time says which counts the trees and the network cannot do
without.
\fi
\item \textbf{Short clocks.} More rapid games, or blitz, until the groups
under a minute are large enough to read.
\item \textbf{The positions that can still be blundered.} The same question
without the lost positions, as a task of its own, scored on its own floor.
\item \textbf{A piece lost or a win not taken.} Two labels where there is
one, so the app can say which kind of mistake a position invites.
\item \textbf{One body for the blunder and the think.} The seconds a move
takes know something about a position that these learners do not. A body
trained on both labels is the run that asks whether it can be read off the
board.
\item \textbf{A blend.} The network's lead is proven early in a game and
not late. An average of the two, or one that knows the stage, may beat
both.
\item \textbf{A stronger body to start from.} The body was trained once, at
one size, on a part of the evaluations file. One trained for longer, on more
of the file or at a larger size has not been tried.
\item \textbf{The game as a sequence.} The moves and the clock use that led
to a position say whether a player is drifting or in time trouble.
\item \textbf{The player's own scale.} A check of the stated chance on the
player's own games, kept private, before the app states it.
\item \textbf{A full exchange search} in place of the count of pieces that
can be taken.
\end{itemize}

\bibliographystyle{unsrt}
\bibliography{../../../papers/references}

\clearpage
\appendix

\section{Sizes, schedules and features}

\begin{table}[htbp]\centering\footnotesize
\block{settings}
\caption{The ladder networks' sizes and schedules, from their run configs.
\emph{Rate} is the peak learning rate. A setting they all share is in the
line under the table.}
\label{tab:settings}
\end{table}

\begin{table}[htbp]\centering\footnotesize
\block{hand}
\caption{The hand features: what \V{task.reference} reads beside the context
numbers, and a network when its run asks for them.}
\label{tab:hand}
\end{table}

\iflook
\section{The positions the model is surest of}

\examples{examples-highest}{The test positions with the highest stated
chance of a blunder.}{fig:highest}

\examples{examples-lowest}{The test positions with the lowest stated chance
of a blunder.}{fig:lowest}

\section{Every cut of the test positions}

Each table below groups the test positions by one column. The columns are the
group, its test rows, the share of moves that were blunders, then for each of
the two runs the chance it stated on average and its PR-AUC, and the paired
gain of the first run over the second with its 95\% interval, bold when the
interval leaves zero out.

\block{slices}
\fi

\end{document}
